\Subhead{5.1.5}{Эффективная площадь и шумовая температура приёмной антенны}

Приёмная эффективная площадь определяется доступной согласованной мощностью $P_{\rm av}=S_{\rm inc}A_e$ для совпадающих поляризаций. Для взаимной антенны
\begin{equation}\label{eq:antennaarea}
 A_e(\theta,\varphi)=\frac{\lambda^2}{4\pi}G(\theta,\varphi).
\end{equation}
Отсюда при свободном пространстве получается формула Фрииса $P_r=P_tG_tG_r[\lambda/(4\pi R)]^2$, если дальняя зона, согласование и поляризация обеспечены. Это формула мощности, не напряжённости.

В радиочастотном классическом пределе доступная шумовая мощность $P_n=k_BT_A\Delta f$. Для идеальной однополяризационной антенны в поле яркостной температуры $T_B(\Omega)$
\begin{equation}
 T_A=\frac1{4\pi}\int D(\Omega)T_B(\Omega)\dd\Omega,
 \qquad \int D\dd\Omega=4\pi.
\end{equation}
Боковые лепестки, видящие тёплую Землю, повышают шум даже при холодном главном луче. При КПД $\eta$ и физической температуре потерь $T_p$ на выходе появляется $\eta T_A+(1-\eta)T_p$. К нему добавляется эквивалентная шумовая температура приёмника. Для сигнала со спектральной плотностью потока, суммированной по двум неполяризованным компонентам, приём одной поляризации даёт дополнительный фактор $1/2$; его нельзя вводить повторно, если $S$ уже относится к принятой поляризации.

Разделив интеграл по главному лепестку и остальным направлениям, получаем
\begin{equation}
 T_A=(1-\beta)\overline T_{\rm г}+\beta\overline T_{\rm вне}.
\end{equation}
Обе температуры усреднены с весом диаграммы по своим областям. Поэтому даже малый коэффициент рассеяния $\beta$ заметен, если боковые лепестки видят тёплую Землю. Это внешний шум идеальной антенны; её потери и шум приёмника добавляются отдельно.
% Шарков, §5.4, с.184, (5.38)–(5.40); обозначения приведены к eq:antennascattering.
